\pgfplotstableread{\sbdata{twists.dat}}\sbTwTab
\pgfplotstablegetrowsof{\sbTwTab}%
\pgfmathtruncatemacro\sbTwLast{\pgfplotsretval-1}%
\gdef\sbTwRow{-1}%
\foreach \i in {0,...,\sbTwLast}{%
  \pgfplotstablegetelem{\i}{q}\of\sbTwTab \let\sbTmpQ\pgfplotsretval
  \pgfplotstablegetelem{\i}{g}\of\sbTwTab \let\sbTmpG\pgfplotsretval
  \ifnum\sbTmpQ=121 \ifnum\sbTmpG=2 \xdef\sbTwRow{\i}\fi\fi}%
\pgfplotstablegetelem{\sbTwRow}{M}\of\sbTwTab \let\sbM\pgfplotsretval
\pgfplotstablegetelem{\sbTwRow}{a}\of\sbTwTab \let\sbA\pgfplotsretval
\pgfplotstablegetelem{\sbTwRow}{b}\of\sbTwTab \let\sbB\pgfplotsretval
\pgfplotstablegetelem{\sbTwRow}{g}\of\sbTwTab \let\sbG\pgfplotsretval
\pgfplotstablegetelem{\sbTwRow}{Nid}\of\sbTwTab \let\sbNid\pgfplotsretval
\pgfmathtruncatemacro\sbCond{\sbM*\sbA+\sbB+1-\sbG}%
\pgfmathtruncatemacro\sbNpl{\sbNid-1}%
\pgfmathtruncatemacro\sbHeur{\sbM*\sbNpl}%
\begin{tikzpicture}[sb, x=0.52cm, y=0.55cm,
    may/.style={draw=sbInk!70, fill=sbOther!30},
    kept/.style={draw=sbInk!70, fill=sbAux},
    zero/.style={draw=sbInk!40, fill=white},
    rowlab/.style={font=\footnotesize, anchor=east, align=right}]
  \def\M{5}
  \node[font=\small, anchor=west] at (-7.4,1.75)
    {(a) the expansion at a place $P$ in a uniformizer $t$, for $M=5$};
  \foreach \j in {0,...,14}
    \node[font=\scriptsize] at (\j+0.5,0.95) {$\j$};
  \node[rowlab] at (-0.2,0.95) {exponent $j$ of $t^j$:};
  \foreach \j in {0,...,14} \draw[may] (\j,-0.4) rectangle ++(1,-0.8);
  \node[rowlab] at (-0.2,-0.8) {a general $z\in\calO_P$};
  \foreach \j in {0,...,14}{%
    \pgfmathtruncatemacro\r{mod(\j,\M)}%
    \ifnum\r=0 \draw[kept] (\j,-1.6) rectangle ++(1,-0.8);
    \else \draw[zero] (\j,-1.6) rectangle ++(1,-0.8);
      \draw[sbInk!40] (\j,-2.4) -- ++(1,0.8);\fi}
  \node[rowlab] at (-0.2,-2.0) {$w^M$};
  \foreach \j in {0,...,14}{%
    \pgfmathtruncatemacro\r{mod(\j,\M)}%
    \ifnum\j>0 \ifnum\r=0 \draw[kept] (\j,-2.8) rectangle ++(1,-0.8);
      \else \draw[zero] (\j,-2.8) rectangle ++(1,-0.8);
        \draw[sbInk!40] (\j,-3.6) -- ++(1,0.8);\fi
    \else \draw[zero] (\j,-2.8) rectangle ++(1,-0.8);
      \draw[sbInk!40] (\j,-3.6) -- ++(1,0.8);
      \draw[sbZero, line width=0.9pt] (\j,-2.8) rectangle ++(1,-0.8);\fi}
  \node[rowlab] at (-0.2,-3.2) {$w^M$ with $w(P)=0$};
  \draw[decorate, decoration={brace, mirror, amplitude=4pt}, sbInk]
    (0,-3.75) -- (5,-3.75)
    node[midway, below=5pt, font=\footnotesize, align=center]
    {$v_P(w^M)\ge M$ from the one condition $w(P)=0$};
  \node[font=\footnotesize, anchor=west, align=left] at (15.4,-2.0)
    {$\Bigl(\sum_j c_jt^j\Bigr)^{M}=\sum_j c_j^{M}t^{jM}$\\[1pt]
     ($M$ a power of $p$)};
  \node[font=\footnotesize, anchor=north west, align=left] at (-7.4,-5.0) {%
    \tikz\draw[may](0,0) rectangle (0.3cm,0.22cm); coefficient may be
nonzero\quad
    \tikz\draw[kept](0,0) rectangle (0.3cm,0.22cm); coefficient $c_j^M$ of
$w^M$\quad
    \tikz{\draw[zero](0,0) rectangle (0.3cm,0.22cm);
      \draw[sbInk!40](0,0)--(0.3cm,0.22cm);} coefficient $0$\quad
    \tikz\draw[zero, draw=sbZero, line width=0.9pt](0,0) rectangle
(0.3cm,0.22cm);
      value at $P$, set to $0$\quad
    \sbschematic[schematic]};
\end{tikzpicture}

\medskip
\begin{tikzpicture}[sb]
  \begin{axis}[sb axis, width=0.8\linewidth, height=0.36\linewidth,
      xmin=0, xmax=150, ymin=0, ymax=1700,
      xlabel={$N$, the number of places where $w^M$ must vanish to order $M$},
      ylabel={linear conditions},
      title={(b) conditions imposed, for $q=121$, $g=2$, $M=\sbM$,
        $(a,b)=(\sbA,\sbB)$},
      title style={font=\small, yshift=8pt},
      clip=false]
    \addplot[sb other plot, sbInk!55, line width=0.9pt, domain=0:150,
      samples=2] {\sbM*x};
    \addplot[sb aux line, line width=1.1pt, domain=0:150, samples=2]
      {\sbCond};
    \draw[sb zero line] (axis cs:\sbNpl,0) -- (axis cs:\sbNpl,1700);
    \node[sb zero] at (axis cs:\sbNpl,\sbHeur) {};
    \node[sb zero] at (axis cs:\sbNpl,\sbCond) {};
    \node[font=\footnotesize, anchor=north west, text=sbInk!70, align=left]
      at (axis cs:4,1640)
      {characteristic-$0$ heuristic: about $MN$ conditions\\
       for independent jets (item (B1)); at $N=\sbNpl$: $\sbHeur$};
    \node[font=\footnotesize, anchor=south west, text=sbInk, align=left,
      inner sep=1pt]
      at (axis cs:40,\sbCond+20)
      {Step~1, characteristic $p$: $\ell(MaP_0+bP_0)=\sbCond$\\
       conditions, for any $N$};
    \node[font=\footnotesize, anchor=south, text=sbZero]
      at (axis cs:\sbNpl,1700) {$N=\sbNpl$};
    \node[font=\footnotesize, anchor=south east, text=sbZero, align=right,
      fill=white, inner sep=1pt, fill opacity=0.85, text opacity=1]
      at (axis cs:\sbNpl-2,560)
      {Frobenius places $P\ne P_0$\\of the computed curve};
  \end{axis}
\end{tikzpicture}
\caption{Multiplicity in characteristic $p$ (Step~3 of
Theorem~\ref{thm:aux-section}). (a) The local expansion at a place $P$, in a
uniformizer $t$, for $M=5$: a general function may have every coefficient
nonzero (grey); an $M$-th power has nonzero coefficients only in degrees
divisible by $M$ (sky blue), because $(\sum_jc_jt^j)^M=\sum_jc_j^Mt^{jM}$
for $M$ a power of $p$; if it also vanishes at $P$ (blue outline), its first
nonzero coefficient is in degree at least $M$. One condition gives order
$M$. Schematic. (b) Linear conditions imposed to make $w^M$ vanish to order
$M$ at $N$ places, for $q=121$, $g=2$, $M=11$, $(a,b)=(10,15)$: in Step~1
the number is $\ell(MaP_0+bP_0)=124$, whatever $N$ is (sky blue, dotted);
for independent jets, as in characteristic~$0$, it would be about $MN$
(grey; a heuristic count, item~(B1) of Section~\ref{sec:dictionary}, not a
theorem). Blue: the $133$ Frobenius places other than $P_0$ of the curve of
Figure~\ref{fig:zero-pole-budget}.}
\label{fig:jets}
